\section{讨论}
\label{sec:discuss}

\subsection{本文的结论有多可靠}

需要区分两类结论。

\textbf{第一类是量化的}：NVLink 与 RDMA 的传输耗时相差 16.7 倍、
聚合方案把 TTFT 降低 13.2 倍、专家共享消除 95\% 重复权重。
这些数字都能追到具体证据与原文（\texttt{analysis.json}），
但它们都是\textbf{别人在别人的条件下测的}。换硬件、换模型、换负载，
数值会变。故本文引用它们时始终带上条件，不给孤立数字。

\textbf{第二类是判据性的}：分离的净收益 ≈ 收益侧的配置优化
$-$ 成本侧的 KV 搬运。这个式子的形式是可辩护的（它只是收支相抵），
但本文\textbf{没有}给出各项的系数——那需要自己的实验。
故第 \ref{sec:boundary} 节的四项判据应被当作\textbf{检查清单}，
而不是可代入计算的模型。

\subsection{一个反复出现的模式}

整理这些材料时，一个模式反复出现：\textbf{几乎所有的收益都来自
「按阶段区分对待」，而不是来自「分离」这个具体手段}。

chunked prefill 区分了两阶段的时间片；分离区分了两阶段的硬件；
逐阶段量化区分了两阶段的精度；EPD 把区分扩展到三段。
它们的共同点是承认「prefill 与 decode 要的东西不同」，
而分离只是实现这个承认的一种方式——且是成本较高的一种
（要搬 KV、要多实例、故障域变大）。

这个视角解释了为什么 chunked prefill 先在 vLLM 里成为默认：
它是「区分对待」的最低成本实现。也解释了为什么分离在有 NVLink 的
同机场景最容易成功：那时它的成本被压到接近零，就退化成了
一种便宜的「区分对待」。

\subsection{没做的事}

本文的局限必须写清楚，否则读者会高估结论。

\begin{itemize}[leftmargin=1.3em,itemsep=1.5pt,topsep=2pt]
  \item \textbf{没有端到端复现。} 未在任何硬件上跑过分离部署，
        所有数字都是文献整理。
  \item \textbf{没有给配比公式。} 第 \ref{sec:split} 节说配比应由
        goodput 优化求解，但本文没有实现这个求解，也没给出
        可代入的参数。
  \item \textbf{模型覆盖不全。} MLA 有 DeepSeek 系、GQA 有 Llama/Qwen 系、
        MoE 有 ExpertPlex，但每类只取了一到两例，
        不足以支撑「某类模型都该这样做」的断言。
  \item \textbf{缺失两篇全文。} 46 篇纳入文献中有 2 篇未取到全文
        （\texttt{2603.04716}、\texttt{2510.08544}），
        它们的主张只能靠其他文献间接支撑，本文未引用其具体数字。
\end{itemize}

\subsection{如果要做下一步实验}

若要把判据变成可计算的模型，最小可行的实验是：在同一套硬件上，
固定模型与 SLO，扫描\textbf{输入输出长度比}这一个变量，
测出「分离 vs chunked prefill」的收益交叉点。
这个交叉点的位置，就是判据 1 与判据 3 的联合定量形式。
只需一套硬件、一个模型、一个 SLO，即可得到一条有实用价值的曲线。

\section{结论}

PD 分离的技术栈可以分成三层：\textbf{为什么分}（两阶段资源画像相反）、
\textbf{怎么搬}（NVLink / RDMA / 存储三档，相差一个数量级以上）、
\textbf{怎么切}（chunked prefill、显式分离、弹性调整三条路线）。

不同模型的最佳实践不能被表述为一套配置，因为最优做法取决于
五个变量（KV 结构、上下文长度、输入输出比、互联拓扑、SLO），
其中只有一个是模型属性。可迁移的结论是判据而非配置：
两阶段最优配置显著不同、存在高带宽 KV 通路、负载足够重、
KV 规模足够大——四项全满足时分离才划算。

最后一点值得重复：文献里报告的加速倍数都是这个判据在特定条件下的
取值。把它当作分离的固有属性，是当前实践中最容易犯的错误。
